\section{模型假设}

（1）几何与声道参数假设。题目给出的等效管道半径、截面积、声道长度、声道位置及固定积分权重均视为测试前已知常数；五条等效声道观测量分别表征不同截面高度的平均流速贡献。

（2）数据使用假设。标准体积只作为训练标签和评价基准；日期、窗口编号及人工扰流编号仅用于离线分组、验证和机理解释，不进入在线预测。对相近流量与剖面状态，短时间内的主要误差可由剖面、动态和零点特征描述。

（3）模型复杂度假设。样本数量有限，候选补偿器采用参数规模受控的线性方法或带正则化的核方法；高维模型仅用于检验样本内表达能力，不将偶然结构解释为未知工况上的稳定规律。
